%
% Chapitre "Structure d'un rapport technique"
%

\chapter{Description}
\label{s:description}

L'équipe d'ingénieurs de NordFaune a été mandatée afin de produire le design conceptuel d'un dispositif autonome fixe destiné à l'échantillonnage de la faune arctique. Cet équipement, conçu pour le Grand Nord, doit permettre l'échantillonnage des activités de lemmings évoluant sur un site arctique de façon complètement passive. Il doit être en mesure d'automatiser et de rendre la mesure autonome pour une période d'un an. 

\bigskip
Une première utilisation consiste à recueillir des images et à collecter les informations à des fins statistiques, tout en assurant une qualité constante des mesures. Ces données devront être archivées pour une période minimale d'au moins un an. La solution proposée vise également à documenter les statistiques sur le territoire tout en réduisant les coûts liés aux relevés terrain. Elle doit assurer une interaction externe minimale en cas de problème. 

\iffalse
\bigskip
L'installation doit permettre une communication hebdomadaire avec une centrale située à plus de 2000 km, offrir une console locale ainsi qu'une application mobile pour l'accès à distance, et fournir un accès sécurisé à la centrale pour trois personnes. Le dispositif devra également être capable de capturer des vidéos visibles ou en proche infrarouge (NIR) d'une durée minimale de 5 secondes, avec une résolution d'au moins 720p à 15 images par seconde. Le temps maximal entre deux vidéos est fixé à une heure. Les vidéos, les paramètres de configuration ainsi que les alarmes devront être conservés pour une période minimale d'un an. Une génération automatique d'alarmes devra être assurée lorsqu'une ou plusieurs fonctionnalités ne sont pas utilisables.
\fi 

\bigskip
Enfin, l'équipement devra être déployé dans un environnement nordique isolé tout en restant fonctionnel dans des conditions de température comprises entre -20 °C et +20 °C. De plus, il doit être fixé par le client, malgré un accès très limité au site. NordFaune doit donc proposer une solution fiable, durable et autonome afin de répondre au besoin de comptage et de documentation de la faune arctique.
